\section{\methodname{} additional details}\label{sec:supp_method}
In this section we provide additional details on TDS.
\Cref{sec:VP_VE_TDS} describes its generalization to alternative common formulations of diffusion models.
\Cref{sec:TDS_hyperparams} describes the choices of proposal variance and resampling strategy.
\Cref{sec:twisted_intermediates_derivation} provides a derivation of \Cref{eq:twisted-intermediates}.
\Cref{sec:supp_method_inpaint_alt_final_step} describes modifications to the final step of TDS in inpainting and inpainting with degrees of freedom applications.
And \Cref{sec:theorem_statement_and_proof} provides a full statement and proof of \Cref{prop:twisted_SMC}.


\subsection{Diffusion models with non-constant noise variance and variance-preserving scaling}\label{sec:VP_VE_TDS}
TDS as developed in \Cref{sec:method} assumed noise with constant variance $\sigma^2$ added at each timestep.
This choice corresponds to a particular discretization of a ``variance exploding'' (VE) diffusion model \citep{song2020score}.
Here we describe generalizations of TDS to non-constant variance schedules,
and to variance preserving (VP) diffusion models \citep{ho2020denoising}.
We adopt the VP formulation in our experiments;
though it introduces additional notation, it is known to work well in practice \citep{song2020score}.

\paragraph{TDS algorithm for variance exploding diffusion models.}
VE diffusion models define the forward process $q$ by 
\begin{align}\label{eq:VE-forward-q}
    q(\dt{x}{1:T} \mid \xzero) := \prod_{t=1}^T q(\xt \mid \xprev), \qquad q(\xt \mid \xprev) := \mathcal{N}(\xt; \xprev, \sigma_t^2). 
\end{align}
where $\sigma_t^2$ is an increasing sequence of variances such that $q(\xT) \approx \mathcal{N}(0,\bar \sigma_T^2)$, with  $\bar \sigma_t^2 :=\sum_{t^\prime=1}^t \sigma_{t^\prime}^2$ for $t=1,\cdots, T$.  And so one can set $p(\xT) =  \mathcal{N}(0,\bar \sigma_T^2)$ to match $q(\xT)$.
Notably, \Cref{eq:VE-forward-q} implies the conditional $q(\xt \mid \xzero) = \mathcal{N}(\xt; \xzero, \bar\sigma_t^2)$. 

The reverse diffusion process $\pmodel$ is parameterized as  
\begin{align}
    \pmodel(\dt{x}{0:T}) :=p(\xT) \prod_{t=T}^{1}\pmodel(\xprev \mid \xt), \qquad 
    \pmodel(\xprev \mid \xt ) := \mathcal{N}(\xprev; \xt + \sigma_t^2 \scoremodel (\xt, t), \sigma_t^2)
\end{align}
where the score network $\scoremodel$ is modeled through a denoiser network $\denoiser$ by $\scoremodel(\xt, t) := (\denoiser (\xt, t; \theta) - \xt )/\bar \sigma_t^2$.
Note that the constant schedule is a special case where $\sigma_t^2=\sigma^2$, and $\bar \sigma_t^2=t\sigma^2$ for all $t$. 

The TDS algorithm for general VE models is described in \Cref{alg:TDS}, where $t\sigma^2$ in \Cref{alg_line:conditional_score} is replaced by $\bar \sigma_t^2$ and $\sigma^2$ in \Cref{alg_line:proposal} is replaced by $\sigma_t^2$. 



\paragraph{Extension to variance preserving diffusion models.}  Another widely used diffusion framework is variance preserving (VP) diffusion models \citep{ho2020denoising}. VP models define the forward process $q$ by 
\begin{align}\label{eq:VP-forward-q}
q(\dt{x}{1:T} \mid \xzero):=\prod_{t=1}^T q(\xt \mid \dt{x}{t-1}), \qquad 
    q(\xt \mid \dt{x}{t-1}) := \mathcal{N}(\xt; \sqrt{1-\sigma_t^2}\dt{x}{t-1}, \sigma_t^2)
\end{align}
where $\sigma_t^2$ is  a sequence of increasing variances chosen such that $q(\xT) \approx \mathcal{N}(0,1)$, and so one can set $p(\xT) = \mathcal{N}(0,1)$. Define $\alpha_t := 1-\sigma_t^2$,  $\bar \alpha_t := \prod_{t^\prime=1}^t\alpha_{t^\prime}$, and $\bar \sigma_t^2:=1-\bar \alpha_t$. 
Then the marginal conditional of \cref{eq:VE-forward-q} is $q(\xt \mid \xzero) = \mathcal{N}(\xt; \sqrt{\bar \alpha_t} \xzero, \bar \sigma_t^2)$. The reverse diffusion process $\pmodel$ is parameterized as  
\begin{align}
    \pmodel(\dt{x}{0:T}) :=p(\xT) \prod_{t=T}^{1}\pmodel(\xprev \mid \xt), \quad 
    \pmodel(\xprev \mid \xt ) := \mathcal{N}(\xprev; \frac{1}{\sqrt{\alpha_t}}\xt + \frac{\sigma_t^2}{\sqrt{\alpha_t}} \scoremodel (\xt, t), \sigma_t^2)
\end{align}
where $\scoremodel$ is now defined through the denoiser $\denoiser$ by $\scoremodel(\xt, t) := (\sqrt{\bar \alpha_t} \denoiser (\xt, t; \theta) - \xt )/\bar \sigma_t^2$.


TDS in \Cref{alg:TDS} extends to VP models as well,
where the conditional score approximation in \Cref{alg_line:conditional_score} 
is changed to  
$\tilde s_k \leftarrow (\sqrt{\bar \alpha_{t+1} }\denoiser(\xtpo_k) - \xtpo_k)/\bar \sigma_{t+1}^2 + \nabla_{\xtpo_k} \log \tilde p_k^{t+1}$, 
and the proposal in \Cref{alg_line:proposal} is changed to 
$\xt_k \sim \pTwist (\cdot \mid \xtpo_k, y) :=\mathcal{N}\left(\frac{1}{\sqrt{\alpha_{t+1}}}\xtpo_k + \frac{\sigma_{t+1}^2}{\sqrt{\alpha_{t+1}}} \tilde s_k, \tilde \sigma^2 \right).$


\subsection{TDS parameters}\label{sec:TDS_hyperparams}
\paragraph{Proposal variance.}
The proposal distribution in  \Cref{alg_line:proposal} of \Cref{alg:TDS} is associated with a variance parameter $\tilde \sigma^2$.
In general, this parameter can be dependent on the time step, i.e.\ replacing $\tilde \sigma^2$ 
 by some $\tilde \sigma_{t+1}^2$ in \Cref{alg_line:proposal}. 
Unless otherwise specified, we set $\tilde \sigma_{t+1}^2:=\mathrm{Var}_{\pmodel}[\dt{x}{t} \mid \xtpo]$ the variance of the unconditional diffusion model (typically learned along with the score network during training). 

\paragraph{Resampling strategy.} The mulinomial resampling strategy in \Cref{alg_line:resampling} of \Cref{alg:TDS} can be replaced by other strategies, see \cite[Chapter 2]{naesseth2019elements} for an overview.
In our experiments, we use the systematic resampling strategy.

In addition, one can consider setting an effective sample size (ESS) threshold (between 0 and 1), and only when the ESS is smaller than this threshold, the resampling step is triggered. 
ESS thresholds for resampling are commonly used to improve efficiency of SMC algorithms \citep[see e.g.][Chapter 2.2.2]{naesseth2019elements},
but for simplicity we use TDS with resampling at every step unless otherwise specified.


\subsection{Derivation of \Cref{eq:twisted-intermediates}}\label{sec:twisted_intermediates_derivation}
%: intermediate twisted targets approach the final target}\label{sec:intermediate_TDS_derivation}
\Cref{eq:twisted-intermediates} illustrated how the extended intermediate targets $\target_t(\dt{x}{0:T})$ provide approximations to the final target $\pmodel(\dt{x}{0:T} \mid y)$ that become increasingly accurate as $t$ approaches 0.
We obtain \Cref{eq:twisted-intermediates} by substituting the proposal distributions in \Cref{eq:diffusion-twisted-proposals} and weights in \Cref{eq:diffusion-twisted-weigthing-functions} into \Cref{eq:final-target} and simplifying.
\begin{align}
   \nu_t(\dt{x}{0:T}) &\propto
   \left[p(\xT)\prod_{t^\prime=0}^{T-1} \pTwist(\dt{x}{t^\prime} \mid \dt{x}{t^\prime+1}, y) \right]
   \left[
  \tilde p(y\mid \xT) \prod_{t^\prime=t}^{T-1} 
    \frac{\pmodel (\dt{x}{t^\prime} \mid \dt{x}{t^\prime +1}) \pTwist(y\mid \dt{x}{t^\prime})}{\pTwist(y \mid \dt{x}{t^\prime+1}) \pTwist(\dt{x}{t^\prime} \mid \dt{x}{t^\prime+1}, y)}
   \right]\\
   & \textrm{Rearrange \ and\ cancel\ }\pmodel \mathrm{\ and \ } \pTwist \textrm{\ terms.}\\ 
   &= \left[ p(\xT) \prod_{t^\prime=t}^{T-1} \pmodel(\dt{x}{t^\prime} \mid  \dt{x}{t^\prime+1}) \right] \left[\tilde p_\theta(y\mid \dt{x}{t})\prod_{t^\prime=0}^{t-1} \pTwist(\dt{x}{t^\prime} \mid \dt{x}{t^\prime +1}, y)\right]\\
   & \textrm{Group\ }\pmodel \textrm{\ terms by chain rule of probability.}\\ 
   &=  \pmodel(\dt{x}{t:T}) \pTwist (y\mid \xt) \prod_{t^\prime=0}^{t-1} \pTwist (\dt{x}{t^\prime} \mid \dt{x}{t^\prime +1}, y)\\
   & \textrm{Apply\ Bayes'\ rule and note that } \pmodel (y\mid \dt{x}{t:T}) = \pmodel(y\mid \xt). \\
   &\propto \pmodel(\dt{x}{t:T}\mid y) \frac{\tilde p_\theta(y\mid \xt)}{p(y \mid \xt)} \prod_{t^\prime=0}^{t-1} \pTwist (\dt{x}{t^\prime} \mid \dt{x}{t^\prime +1}, y)\\
%   & \textrm{Multiply\ by\  }\prod_{t^\prime=0}^{t-1} \pmodel(\dt{x}{t^\prime}\mid \dt{x}{t^\prime + 1}, y)/ \pmodel(\dt{x}{t^\prime}\mid \dt{x}{t^\prime + 1}, y)=1.\\
 & \textrm{Note that }  \pmodel(\dt{x}{0:T} \mid y) = \pmodel (\dt{x}{t:T} \mid y)   \prod_{t=0}^{t'-1} \pmodel(\dt{x}{t^\prime}\mid \dt{x}{t^\prime + 1}, y). \\
   &=  \pmodel(\dt{x}{0:T}\mid y) \left[\frac{\tilde p_\theta(y\mid \xt)}{p(y \mid \xt)} \prod_{t^\prime=0}^{t-1} \frac{ \pTwist (\dt{x}{t^\prime} \mid \dt{x}{t^\prime +1}, y)}{
   \pmodel(\dt{x}{t^\prime}\mid \dt{x}{t^\prime + 1}, y)}\right].
\end{align}
The final line is the desired expression in \Cref{eq:twisted-intermediates}.
\subsection{Inpainting and degrees of freedom final steps for asymptotically exact target }\label{sec:supp_method_inpaint_alt_final_step}
The final ($t=0$) twisting function for inpainting and inpainting with degrees of freedom described in \Cref{subsec:conditioning-problems} do not satisfy the assumption of \Cref{thm:twisted_SMC_complete} that $\pTwist (y\mid \xzero)=\plikelihood{y}{\xzero}.$
This choice introduces error in the final target of TDS relative to the exact conditional $\pmodel(\dt{x}{0:T} \mid y).$

For inpainting, to maintain asymptotic exactness one may instead choose the final proposal and weights as 
\ifnum\workshopsubmission=0
\begin{align}\label{eqn:inpaint_final_proposal_and_weight}
    \pTwist(\xzero \mid \dt{x}{1}, y; \mask) :=  \delta_{y}(\xzero_\mask)\pmodel(\xzero_\umask \mid \dt{x}{1}) \quad \mathrm{and} \quad \weight_0(\xzero, \dt{x}{1}) = 1.
\end{align}
\else
\begin{align}\label{eqn:inpaint_final_proposal_and_weight}
    \pTwist(\xzero \mid \dt{x}{1}, y; \mask) :=  \delta_{y}(\xzero_\mask)\pmodel(\xzero_\umask \mid \dt{x}{1})
\end{align}
and $\weight_0(\xzero, \dt{x}{1}) = 1.$
\fi
One can verify the resulting final target is $\pmodel(\xzero \mid \xzero_\mask =y)$ according to \cref{eq:final-target}.  

Similarly, for inpainting with degrees of freedom, one may define the final proposal and weight as
\ifnum\workshopsubmission=0
\begin{equation}
    \pTwist(\xzero \mid \dt{x}{1}, y, \mathcal{M}) :=  \sum_{\mask \in \mathcal{M}} \frac{\pmodel(\xzero_\mask \mid \dt{x}{1}, \mask)}{\sum_{\mask^\prime \in \mathcal{M}} \pmodel(\xzero_{\mask^\prime} \mid \dt{x}{1}, \mask^\prime)
    }
    \delta_{y}(\xzero_\mask)\pmodel(\xzero_\umask \mid \dt{x}{1}) 
    \quad \mathrm{and} \quad \weight_0(\xzero, \dt{x}{1}) = 1.
\end{equation}
\else
$\weight_0(\xzero, \dt{x}{1}) = 1$ and 
\begin{align*}
\begin{split}
    &\pTwist(\xzero \mid \dt{x}{1}, y; \mathcal{M}) :=  \sum_{\mask \in \mathcal{M}} \pmodel(\xzero_\mask \mid \dt{x}{1};\mask)\delta_{y}(\xzero_\mask)\pmodel(\xzero_\umask \mid \dt{x}{1}).
\end{split}
\end{align*}
\fi

\subsection{Asymptotic accuracy of TDS -- additional details and full theorem statement}\label{sec:theorem_statement_and_proof}
In this section we (i) characterize sufficient conditions on the model and twisting functions under which TDS provides arbitrarily accurate estimates as the number of particles is increased and 
(ii) discuss when these conditions will hold in practice for the twisting function $\pTwist(y\mid \xt)$ introduced in \Cref{sec:method}.

We begin with a theorem providing sufficient conditions for asymptotic accuracy of TDS.
\begin{theorem}\label{thm:twisted_SMC_complete}
Let $\pmodel(\dt{x}{0:T})$ be a diffusion generative model
(defined by \cref{eqn:diffusion_marginal,eq:unconditional-diffusion-transition})
with 
$$
\pmodel(\xt \mid \xtpo) = \mathcal{N}\left(\xt \mid \xtpo + \sigma_{t+1}^2 \scoremodel(\xtpo), \sigma_{t+1}^2 I
\right),$$
with variances $\sigma_1^2, \dots, \sigma_T^2.$ 
Let $\pTwist(y \mid \xt)$ be twisting functions,
and 
$$
\proposal_t(\xt \mid \xtpo) = 
\mathcal{N}\left(
\xt \mid \xtpo + \sigma_{t+1}^2 [\scoremodel(\xtpo) + \nabla_{\xtpo}\log \pTwist(y\mid \xtpo)], 
\tilde \sigma_{t+1}^2 \right)
$$
be proposals distributions for $t=0,\dots, T-1,$
and let $\mathbb{P}_{K}=\sum_{k=1}^K \dt{w_k}{0}\delta_{\dt{x_k}{0}}$ for weighted particles $\{(\dt{x_k}{0}, \dt{w_k}{0})\}_{k=1}^K$ returned by \Cref{alg:TDS} with $K$ particles.
Assume
\begin{enumerate}[label=(\alph*)]
    \item{the final twisting function is the likelihood, $\pTwist(y\mid \xzero)=\plikelihood{y}{\xzero},$} \label{item:final_twist}
    \item{the first twisting function $\pTwist(y \mid \xT),$ and the ratios of subsequent twisting functions $\pTwist(y \mid \xt) / \pTwist(y \mid \xtpo) $ are positive and bounded,}\label{item:twist_assumption}
    %\item{each $\nabla_{\xt} \log \pTwist(y\mid \xt)$ with $t>0$ is continuous and has bounded gradients, and}\label{item:twist_gradient}
    \item{each $\log \pTwist(y\mid \xt)$ with $t>0$ is continuous and has bounded gradients in $\xt$, and}\label{item:twist_gradient}
    \item{the proposal variances are larger than the model variances, i.e.\ for each $t,\ \tilde \sigma_t^2 >\sigma_t^2.$} \label{item:proposal_var}
\end{enumerate}
Then $\mathbb{P}_{K}$ converges setwise to $\pmodel(\xzero\mid y)$ with probability one,
that is 
for every set $A,$ 
$\lim_{K\rightarrow\infty} \mathbb{P}_{K}(A) = \int_A \pmodel(\xzero\mid y) d\xzero.$
\end{theorem}

The assumptions of \Cref{thm:twisted_SMC_complete} are readily satisfied in common applications.
\begin{itemize}
\item {Assumption \ref{item:final_twist} may be satisfied by construction by choosing $\pTwist(y\mid \xzero)=\plikelihood{y}{\xzero}$
by defining $\denoiser(\xzero, t=0)=\xzero.$
}
\item{Assumption \ref{item:twist_assumption} is satisfied if (i) $\plikelihood{y}{x}$ is smooth in $x$ and everywhere positive and (ii) $\pTwist(y\mid \xt) = \plikelihood{y}{\denoiser(\xt)}$ where $\denoiser(\xt)$ takes values in some compact domain.
An alternative sufficient condition for Assumption \ref{item:twist_assumption} is for $\plikelihood{y}{x}$ to be positive and bounded away from zero; 
this latter condition will hold when, for example, $\plikelihood{y}{x}$ is a classifier fit with regularization.}
\item{Assumption \ref{item:twist_gradient} is the strongest assumption.
It will be satisfied, for example, if (i) $\pTwist(y\mid \xt) = \plikelihood{y}{\denoiser(\xt)},$ and (ii) $\plikelihood{y}{x}$ and $\denoiser(\xt)$ are smooth in $x$ and $\xt,$ with uniformly bounded gradients.
While smoothness of $\denoiser(\cdot, t)$ can be encouraged by the use of skip-connections and regularization,
$\denoiser(\cdot, t)$ may present sharp transitions,
particularly for $t$ close to zero.}
\item{Assumption \ref{item:proposal_var}, that the proposal variances satisfy $\tilde \sigma_{t}^2 > \sigma_t^2$ is likely not needed for the result to hold.
However, this assumption permits usage of existing SMC theoretical results in the proof;
in practice, our experiments use $\tilde \sigma_t^2 = \sigma_t^2,$
but alternatively the assumption could be met by inflating each $\tilde \sigma_t^2$ by some arbitrarily small $\delta$ without markedly impacting the behavior of the sampler.}
\end{itemize}

%\paragraph{\BLT{Cut this paragraph for NeurIPS appendix submission} Inpainting and inpaiting with degrees of freedom.}

%\begin{cor}
%The inpainting and DOF cases are consistent for their conditionals.
%\end{cor}
%\begin{proof}
%We can see that early truncation at step $t=1$ forms a special case.
%By reindexing $t$ we find we are consistent for $\pmodel(\dt{x}{1} \mid y).$
%Then those final samplers sample $\xzero \sim \pmodel(\xzero \mid \dt{x}{1}, y)$ in the last step.
%So by the law of total probability and chain rule of probability they give exact samples in the limit as desired.
%\end{proof}
%
%Why is bounded assumption reasonable in the inpainting case?
%It's because each $\sigma_t^2>0$ and so with 
%$\pTwist(y\mid \xt; \mask) = \mathcal{N}(y; \denoiser(\xt)_\mask, \bar \sigma_t^2)=p(y\mid \dt{x}{1})$
%with $\bar \sigma_t^2 = \sum_{t^\prime=1}^t\sigma_{t^\prime}^2,$
%the twisting variance are a decreasing sequence.
%And since the denoiser is also bounded and each $\bar \sigma_t^2>0$, there is some $0 <c <C< \infty$ such that 
%$c<\pTwist(y\mid \xt) <C$ and so the ratio is bounded.

\paragraph{Proof of \Cref{thm:twisted_SMC_complete}:}
\Cref{thm:twisted_SMC_complete} characterizes a set of conditions under which SMC algorithms converge.
We restate this result below in our own notation.
\begin{theorem}[\citet{chopin2020introduction} -- Proposition 11.4]\label{thm:chopin_theorem}
Let $\{(\dt{x_k}{0}, \dt{w_k}{0})\}_{k=1}^K$ be the particles and weights returned at the last iteration of a sequential Monte Carlo algorithm with $K$ particles using multinomial resampling.
If each weighting function $\weight_t(\xt, \xtpo)$ is positive and bounded, 
then for every bounded, $\target_0$-measurable function $\phi$ of $\xt$
$$
\lim_{K\rightarrow \infty} \sum_{k=1}^K \dt{w_k}{0} \phi(\dt{x_k}{0})
=
\int \phi(\xzero)\target_0(\xzero)d\xzero.
$$
with probability one.
\end{theorem}


An immediate consequence of \Cref{thm:chopin_theorem} is the setwise convergence of the discrete measures, $\hat P_K = \sum_{k=1}^K \dt{w_k}{0} \delta_{\dt{x_k}{0}}.$
This can be seen by taking for each $\phi(x) = \I[x \in A]$ for any $\target_0$-measurable set $A$.
The theorem applies both in the Euclidean setting, where each $\dt{x_k}{t}\in \R^D,$ as well as the Riemannian setting.

We now proceed to prove \Cref{thm:twisted_SMC_complete}.
\begin{proof}
To prove the theorem we show
(i) the $\xzero$ marginal final target $\target_0$ is $\pmodel(\xzero \mid y)$ and then
(ii) $\mathbb{P}_K$ converges setwise to $\target_0.$

We first show (i) by manipulating $\target_0$ in \Cref{eq:final-target} to obtain $\pmodel(\dt{x}{0:T} \mid y).$
From \Cref{eq:final-target} we first have
\begin{align}
\begin{split}
   \target_0(\dt{x}{0:T}) &= \frac{1}{\mathcal{L}_0} 
        \left[ \proposal(\dt{x}{T})  \prod_{t=0}^{T-1} \proposal_{t}(\dt{x}{t}\mid  \dt{x}{t+1})  \right] 
        \left[ \weight_T(\dt{x}{T})
        \prod_{t=0}^{T-1} \weight_{t^\prime}(\dt{x}{t}, \dt{x}{t+1}) \right]   \\
    &\mathrm{Substitute\ in\ weights\ from\ \cref{eq:diffusion-twisted-weigthing-functions}.} \\
    &= \frac{1}{\mathcal{L}_0} 
        \left[ p(\dt{x}{T}) \prod_{t=0}^{T-1} r_t (\xt \mid \xtpo)
        \right] 
        \left[ \pTwist(y\mid \dt{x}{T})
        \prod_{t=0}^{T-1} 
        \frac{\pmodel(\xt \mid \xtpo)\pTwist(y\mid \xt)}{\pTwist(y \mid \xtpo)r_t(\xt \mid \xtpo)}
        \right] \\
    &\mathrm{Rearrange\ }\pmodel\mathrm{\ and \ } r_t\ \mathrm{\ terms, \ and\ cancel\ out\ }r_t \mathrm{\ terms.} \\
    &= \frac{1}{\mathcal{L}_0} 
        \left[ p(\dt{x}{T}) \prod_{t=0}^{T-1} \pmodel(\xt \mid \xtpo)
        \right] 
        \left[ \pTwist(y\mid \dt{x}{T})
        \prod_{t =0}^{T-1} 
        \frac{\cancel{r_t(\xt \mid \xtpo)}\pTwist(y\mid \xt)}{\pTwist(y \mid \xtpo)\cancel{r_t(\xt \mid \xtpo)}}
        \right] \\
    &\mathrm{Collapse\ }\pmodel\mathrm{\ terms.}\\ %,\ rearrange\ and\ cancel\ }\pTwist \mathrm{\ terms.} \\
    &= \frac{1}{\mathcal{L}_0} 
        \pmodel(\dt{x}{0:T})
        \left[ 
        \prod_{t=0}^{T-1} 
        \frac{\pTwist(y\mid \xt)}{\pTwist(y \mid \xtpo)}\right]\pTwist(y \mid \dt{x}{T})\\
    &\mathrm{Cancel \ out \ } \pTwist \mathrm{\ terms.}\\
    &= \frac{1}{\mathcal{L}_0} 
        \pmodel(\dt{x}{0:T}) \pTwist (y \mid \xzero) \\ 
    &\mathrm{Recognize\ }\pTwist(y\mid \dt{x}{0})=\plikelihood{y}{\xzero}\mathrm{\ by \ Assm. \ref{item:final_twist}\ \ and\  apply\  Bayes'\  rule\  with\  }\mathcal{L}_0=\pmodel(y). \\
    &= \pmodel(\dt{x}{0:T}\mid y).
\end{split}
\end{align}
The final line reveals that once we marginalize out $\dt{x}{1:T}$
we obtain $\target_0(\xzero)=p(\xzero \mid y)$ as desired.

We next show that $\mathbb{P}_K$ converges to $\target_0$ with probability one by applying \Cref{thm:chopin_theorem}.
To apply \Cref{thm:chopin_theorem} it is sufficient to show that the weights at each step are upper bounded, as they are defined through (ratios of) probabilities and hence are positive. 
Since there are a finite number of steps $T,$
it is enough to show that each $w_t$ is bounded.
The inital weight is the initial twisting function, $w_T(\xT)=\pTwist(y\mid \xT),$ which is bounded by Assumption \ref{item:twist_assumption}.
So we proceed to intermediate weights.

To show that the weighting functions at subsequent steps are bounded,
we decompose the log-weighting functions as 
$$
\log \weight_t(\xt, \xtpo) = \log \frac{\pTwist(y\mid \xt)}{\pTwist(y\mid \xtpo)} + \log \frac{\pmodel(\xt\mid \xtpo)}{r_t(\xt \mid \xtpo)} ,
$$
and show independently that $\log \pTwist(y\mid \xt)/\pTwist(y\mid \xtpo)$ and
$\log \pmodel(\xt\mid \xtpo)/r_t(\xt \mid \xtpo)$ are bounded.
The first term $\log \pTwist(y\mid \xt)/\pTwist(y\mid \xtpo)$ is again bounded by Assumption \ref{item:twist_assumption},
and we proceed to the second.

That $\log \pmodel(\xt\mid \xtpo)/r_t(\xt \mid \xtpo)$ is bounded follows from Assumptions \ref{item:twist_gradient} and \ref{item:proposal_var}.
First write
$$
%\pmodel(\xt \mid \xtpo) = \mathcal{N}\left(\xt \mid \xtpo + \sigma_{t+1}^2 \scoremodel(\xtpo), \sigma_{t+1}^2 I
\pmodel(\xt \mid \xtpo) = \mathcal{N}\left(\xt \mid \hat \mu, \sigma_{t+1}^2 I
\right)$$
with $\hat \mu=\xtpo + \sigma_{t+1}^2 \scoremodel(\xtpo),$
and 
$$
\proposal_t(\xt \mid \xtpo) = 
\mathcal{N}\left(
\xt \mid  \hat \mu_\psi, \tilde \sigma_{t+1}^2 I \right),
$$
for $\hat \mu_\psi = \hat \mu + \sigma_{t+1}^2 \nabla_{\xtpo}\log \pTwist(y\mid \xtpo)$.
The log-ratio then simplifies as 
\begin{align}
\log \frac{\pmodel(\xt\mid \xtpo)}{r_t(\xt \mid \xtpo)} &= 
\log \frac{| 2\pi \sigma_{t+1}^2 I|^{-1/2} \exp\{-(2\sigma_{t+1}^2)^{-1}\|\hat \mu - \xt \|^2 \}}
{ | 2\pi \tilde \sigma_{t+1}^2 I|^{-1/2} \exp\{-(2\tilde \sigma_{t+1}^2)^{-1}\|\hat \mu_\psi- \xt \|^2 \}} \\
&\textrm{Rearrange \ and\  let  \ } C=\log | 2\pi \sigma_{t+1}^2 I|^{-1/2} /  | 2\pi \tilde \sigma_{t+1}^2 I|^{-1/2}\\
&= \frac{-1}{2}\left[ 
\sigma_{t+1}^{-2}\|\hat \mu - \xt \|^2    -
\tilde \sigma_{t+1}^{-2}\|\hat \mu_\psi - \xt \|^2 \right]  +C \\
&\textrm{Expand and rearrange \ } \|\hat \mu_\psi - \xt \|^2 = \| \hat \mu  - \xt \|^2 + 2 \langle \hat\mu_\psi - \hat \mu, \hat \mu - \xt  \rangle + \|\hat \mu_\psi - \hat \mu\|^2 \\
&= \frac{-1}{2}\left[ ( \sigma_{t+1}^{-2} - \tilde \sigma_{t+1}^{-2})\|\hat \mu - \xt \|^2  
- 2 \tilde \sigma_{t+1}^{-2}\langle \hat\mu_\psi - \hat \mu, \hat \mu - \xt  \rangle - \tilde \sigma_{t+1}^2 \|\hat \mu_\psi - \hat \mu\|^2 
\right]+ C\\
&\textrm{Let \ } C^\prime {=} C - \frac{1}{2}\tilde \sigma_{t+1}^2 \|\hat \mu_\psi - \hat \mu\|^2 \textrm{\ and rearrange.  Note that $\|\hat \mu_\psi {-} \hat \mu\|^2{<}\infty$ by Assm. \ref{item:twist_gradient}.}\\
&= \frac{-1}{2}(\sigma_{t+1}^{-2} - \tilde \sigma_{t+1}^{-2})\|\hat \mu - \xt \|^2 +
\tilde \sigma_{t+1}^{-2}\langle \hat\mu_\psi - \hat \mu, \hat \mu - \xt  \rangle + C^\prime\\
&\textrm{Apply Cauchy-Schwarz}\\
&\le  \frac{-1}{2} (\sigma_{t+1}^{-2} - \tilde \sigma_{t+1}^{-2})\|\hat \mu - \xt \|^2
+ \tilde \sigma_{t+1}^{-2}\| \hat\mu_\psi - \hat \mu\|\cdot \| \hat \mu - \xt  \| + C^\prime\\
&\textrm{Upper-bounding using that $\mathrm{max}_{x}\frac{-a}{2}x^2 + bx =\frac{b^2}{2a}$ for $a = \sigma_{t+1}^{-2} - \tilde \sigma_{t+1}^{-2}>0$ by Assm. (d).}\\
&\le  \frac{1}{2} \frac{(\tilde \sigma_{t+1}^{-2}\| \hat\mu_\psi - \hat \mu\|)^2}{\sigma_{t+1}^{-2} - \tilde \sigma_{t+1}^{-2}} + C^\prime \\
&= \frac{\tilde \sigma_{t+1}^{-4}}{2( \sigma_{t+1}^{-2} - \tilde \sigma_{t+1}^{-2})}\| \hat\mu_\psi - \hat \mu\|^2  + C^\prime \\
& \textrm{Note that } \hat \mu_\psi = \hat \mu + \sigma_{t+1}^2 \nabla_{\xtpo}\log \pTwist(y\mid \xtpo). \\ 
&= \frac{\tilde \sigma_{t+1}^{-4}}{2( \sigma_{t+1}^{-2} - \tilde \sigma_{t+1}^{-2})}\sigma_{t+1}^4 \| \nabla_{\xtpo} \log \pTwist(y \mid \xtpo) \|^2  + C^\prime \\
&\le C^{\prime\prime}.
\end{align}
The final line follows from Assumption \ref{item:twist_gradient}, that the gradients of the twisting functions are bounded.
The above derivation therefore provides that each $\weight_t$ is bounded, concluding the proof.
\end{proof}

\section{Riemannian \methodname{}}\label{sec:riemannian}

This section provides additional details on the extension of TDS to Riemannian diffusion models introduced in \Cref{sec:method}.
We first introduce the tangent normal distribution.
We then provide with background on Riemannian diffusion models, which we parameterize with the tangent normal.
Then we describe the extension of TDS to these models.
Finally we show how \Cref{alg:TDS} modifies to this setting.

\paragraph{The tangent normal distribution.}
Just as the Gaussian is the parametric family underlying Euclidean diffusion models in \Cref{eq:unconditional-diffusion-transition}, 
the tangent normal (see e.g.\ \citep{chirikjian2014gaussian,deriemannian})
underlies generation in Riemannian diffusion models so we review it here.

We take the tangent normal to be the distribution is implied by a two step procedure.
Given a variable $x$ in the manifold, the first step is to sample a variable $\bar y$ in $\mathcal{T}_{x},$ the tangent space at $x;$
if $\{h_1, \dots, h_D\}$ is an orthonormal basis of $\mathcal{T}_{x}$
one may generate
\begin{align}
    \bar y = \mu + \sum_{d=1}^D \sigma \epsilon_d \cdot h_d,
\end{align}
with $\epsilon_d\overset{i.i.d.}{\sim}\mathcal{N}(0, 1),$
and $\sigma^2>0$ a variance parameter. 
The second step is to project $\bar y$ back onto the manifold to obtain
$y = \exp_{x}\{ \bar y\}$
where $\exp_{x}\{ \cdot \}$ denotes the exponential map at $x.$ The resulting distribution of $y$ is denoted as $\mathcal{TN}_{x}(\mu, \sigma^2)$. 
By construction, $\mathcal{T}_x$ is a Euclidean space with its origin at $x,$
so when $\|\mu\|_2=0, \mathcal{TN}_{x}(\mu, \sigma^2)$ is centered on $x.$
And since the geometry of a Riemannian manifold is locally Euclidean, when $\sigma^2$ is small
the exponential map is close to the identity map and the tangent normal is essentially a narrow Gaussian distribution in the manifold at $x.$
Finally, we use $\mathcal{TN}_{x}(y;\mu, \sigma^2)$ to denote the density of the tangent normal evaluated at $y$.

Because this procedure involves a change of variables from $\bar y$ to $y$ (via the exponential map),
to compute the tangent normal density one computes 
\begin{equation}\label{eq:TN_density}
    \mathcal{TN}_{x}(y;\mu, \sigma^2) = \mathcal{N}(\exp^{-1}_x\{y\}; \mu,  \sigma^2) \left| \frac{\partial}{\partial y} \exp_x^{-1}\{y\}\right|
\end{equation}
where $\exp_x^{-1}\{y\}$ is the inverse of the exponential map (from the manifold into $\mathcal{T}_x$), 
and $\left| \frac{\partial}{\partial y} \exp_x^{-1}\{y\}\right|$ is the determinant of the Jacobian of the exponential map; when the manifold lives in a higher dimensional subset of $R,$ we take 
$\left| \frac{\partial}{\partial y} \exp_x^{-1}\{y\}\right|$ to be the product of the positive singular values of the Jacobian.

\paragraph{Riemannian diffusion models and the tangent normal distribution.}
Riemannian diffusion models proceeds through a geodesic random walk \citep{deriemannian}.
At each step $t,$
one first samples a variable $\dt{\bar x}{t}$ in $\mathcal{T}_{\xtpo};$
if $\{h_1, \dots, h_D\}$ is an orthonormal basis of $\mathcal{T}_{\xtpo}$
one may generate
\begin{align}
    \dt{\bar x}{t}  =\sigma_{t+1}^2 s_\theta(\dt{x}{t+1})+ \sum_{d=1}^D \sigma_{t+1} \epsilon_d \cdot h_d,
\end{align}
with $\epsilon_d\overset{i.i.d.}{\sim}\mathcal{N}(0, 1).$
One then projects $\dt{\bar x}{t}$ back onto the manifold to obtain
$\xt = \exp_{\xtpo}\{ \dt{\bar x}{t}\}$
where $\exp_{x}\{ \cdot \}$ denotes the exponential map at $x.$

This is equivalent to sampling $\xt$ from a tangent normal distribution as 
\begin{equation}\label{eqn:reverse_transition_unconditional_riemannian}
    \dt{x}{t} \sim p(\dt{ x}{t} \mid \xtpo) = \mathcal{TN}_{\xtpo}\left(\dt{ x}{t};\sigma_{t+1}^2 s_\theta(\dt{x}{t+1}),
\sigma_{t+1}^2 \right).
\end{equation}

\paragraph{TDS for Riemannian diffusion models.}
To extend TDS, appropriate analogues of the twisted proposals and weights are all that is needed.
For this extension we require that the diffusion model is also associated with a manifold-valued denoising estimate $\denoiser$ as will be the case when, for example, $\scoremodel(\xt, t) := \nabla_{\xt} \log q(\xt \mid \xzero=\denoiser)$ for $\denoiser = \denoiser(\xt, t).$
In contrast to the Euclidean case, a relationship between a denoising estimate and a computationally tractable score approximation may not always exist for arbitrary Riemannian manifolds;
however for Lie groups when the the forward diffusion is the Brownian motion, tractable score approximations do exist \citep[Proposition 3.2]{yim2023se}.

For the case of positive and differentiable $\plikelihood{y}{\xzero},$ we again choose twisting functions $\twist(y\mid \xt):= \plikelihood{y}{\denoiser(\xt)}.$

Next are the inpainting and inpainting with degrees of freedom cases.
Here, assume that $\xzero$ lives on a multidimensional manifold (e.g.\ $SE(3)^N$)
and the unmasked observation $y=\xzero_\mask$ with $\mask \subset \{1, \dots, N\}$ on a lower-dimensional sub-manifold (e.g.\ $SE(3)^{|\mask|},$ with $|\mask|<N$).
In this case, twisting functions are constructed exactly as in \Cref{subsec:conditioning-problems}, 
except with the normal density in \Cref{eqn:inpaint_twist} replaced with a Tangent normal as
\begin{align}
\twist(y\mid \xt) {=} \mathcal{TN}_{\denoiser(\xt)_\mask}( y ; 0, \bar \sigma_t^2).
\end{align}

For all cases, we propose the twisted proposal as
\begin{equation}\label{eqn:likelihood_approx_riemannian_supp}
    \pTwist(\dt{ x}{t} \mid \xtpo, y) {=} \mathcal{TN}_{\xtpo}\left(\dt{ x}{t};\sigma_{t+1}^2 s_\theta(\dt{x}{t+1}, y), \tilde\sigma_{t+1}^2 \right)
\end{equation}
%where $\bar y := \exp_{\denoiser(\dt{x}{t})_\umask}^{-1}(y)\in \mathcal{T}_{\denoiser(\dt{x}{t})_\umask}$ 
%is $\dt{x}{0}$ mapped into $\mathcal{T}_{\denoiser(\dt{x}{t})}$ by the inverse of the exponential map at $\denoiser(\dt{x}{t})_\umask,$
where as in the Euclidean case  $s_\theta(\dt{x}{t}, y)=s_\theta(\dt{x}{t}) + \nabla_{\dt{x}{t}} \log \tilde p(y \mid \dt{x}{t})$.


Weights at intermediate steps are computed as in the Euclidean case (\Cref{eq:diffusion-twisted-weigthing-functions}): 
%However, 
%since the tangent normal includes a projection step back to the manifold, its density involves not only usual Gaussian density but also the Jacobian of the exponential map to account for the change of variables.
%For example, in the case that the ambient dimension of the manifold and the tangent space are the same dimension,
%$$
%\mathcal{TN}_{x}(x^\prime;\mu, \sigma^2)= \mathcal{N}(\exp^{-1}_x\{x^\prime\}; \mu, \sigma^2)\left|\frac{\partial}{\partial x} \exp_x^{-1}\{x^\prime\}\right|,
%$$
%where $\exp^{-1}_x\{x^\prime\}$ is the inverse of the exponential map at $x$ from the manifold into $\mathcal{T}_x,$
%and 
%$\left|\frac{\partial}{\partial x} \exp_x^{-1}\{x^\prime\}\right|$ is the determinant of its Jacobian.
%The weights are then computed as
\begin{align}\label{eq:diffusion-twisted-weigthing-functions-riemannian}
    \weight_t (\xt , \xtpo ):&= \frac{\pmodel (\xt \mid \xtpo) \twist(y \mid \xt)}{\twist(y \mid \xtpo) \pTwist(\xt \mid \xtpo, y)} \\
    &= \frac{\mathcal{TN}_{\xtpo}(\xt; \sigma_{t+1}^2\scoremodel(\xtpo), \sigma_{t+1}^2) \twist(y \mid \xt)
    }{\twist(y \mid \xtpo) \mathcal{TN}_{\xtpo}(\xt; \sigma_{t+1}^2\scoremodel(\xtpo, y), \tilde\sigma_{t+1}^2)}.
\end{align}
While the proposal and target contribute identical Jacobian determinant terms that cancel out, they remain in the twisting functions.


\paragraph{Adapting the TDS algorithm to the Riemannian setting.}
To translate the TDS algorithm to the Riemannian setting we require only two changes.

The first is on \Cref{alg:TDS} 
\Cref{alg_line:conditional_score}.
Here we assume that the unconditional score is computed through the transition density function: 
\begin{align}
\scoremodel(\xtpo) := \nabla_{\xtpo} \log q_{t+1{|}0}(\xtpo\mid \denoiser) \quad \mathrm{\ for\ } \denoiser = \denoiser(\xtpo).
\end{align}
Note that the gradient above ignores the dependence of $\denoiser$ on $\xtpo$. 

%\BLT{Is it clear that the gradient above ignores the dependence of $\denoiser$ on $\xt$?}
The conditional score approximation in \Cref{alg_line:conditional_score} is then replaced with 
$$
\tilde s_k \leftarrow \pTwist (\cdot \mid \xtpo_k, y):= \nabla_{\xtpo_k} \log q_{t+1{|} 0}(\xtpo_k\mid \denoiser) + \nabla_{\xtpo_k} \log \tilde p_k^{t+1} \quad \mathrm{ \ for\  } \denoiser = \denoiser(\xtpo_k).
$$
Notably, $\tilde s_k$ is a $\mathcal{T}_{\xtpo_k}$-valued Riemannian gradient.

The second change is to make the proposal on \Cref{alg:TDS} \Cref{alg_line:proposal} a tangent normal, as defined in \cref{eqn:likelihood_approx_riemannian_supp}.

